\input{preamble.tex}
\title{Free products and recovery of the input geometries}
\author{Human Compression Project}\date{Private draft, 11 October 2026}
\begin{document}\maketitle
We use the convention $A_J=\bigcap_{j\in J}A_j$, with
$A_\varnothing=A$. All type sets in this note are finite. Neither
faithfulness nor a trivial chamber stabilizer is implicit. A family of
subsets is called Helly here if every finite pairwise intersecting
subfamily has nonempty total intersection.

\begin{lemma}\label{lem:cosets}
A coset incidence system $(A,(A_i)_{i\in I})$ is flag-transitive if
and only if the family of all left cosets of all $A_J$, $J\subseteq I$,
is Helly.
\end{lemma}
\begin{proof}
A flag is a collection of pairwise intersecting cosets of distinct
maximal parabolics. It is a translate of a standard flag precisely
when these cosets have a common element. This proves one implication.
For the other, expand each coset $aA_J$ into the intersection of
$aA_j$, $j\in J$. If two original cosets intersect, every maximal
coset from one meets every maximal coset from the other. Within one
expanded coset they meet at $a$. Two such cosets of the same maximal
parabolic are therefore equal. The resulting distinct-type collection
is a flag, and flag transitivity supplies a common element. Empty
$J$ imposes no constraint. A standard flag extends to the standard
chamber, so flag transitivity also supplies the geometric property.
\end{proof}

\begin{lemma}\label{lem:free-helly}
Let $\mathcal P$ and $\mathcal Q$ be families of subgroups of $A$ and
$B$, respectively, such that their left cosets are Helly. Then the
left cosets of the subgroups $P*Q\leq A*B$, with
$P\in\mathcal P$ and $Q\in\mathcal Q$, are Helly. Moreover,
\[
 \bigcap_{j=1}^m(P_j*Q_j)
   =\left(\bigcap_{j=1}^mP_j\right)*
     \left(\bigcap_{j=1}^mQ_j\right).
\]
\end{lemma}
\begin{proof}
The intersection assertion follows from the unique reduced word in a
free product: a word belongs to $P*Q$ exactly when each of its
$A$ syllables belongs to $P$ and each of its $B$ syllables belongs to
$Q$. This also proves $(P*Q)\cap A=P$ and $(P*Q)\cap B=Q$.

Consider the tree with vertices $gA,gB$ and with an edge labelled $g$
joining them, for every $g\in A*B$. For $H=P*Q$, the union of the
edges labelled by elements of $H$ is a connected subtree. Indeed a
reduced word in $P$ and $Q$ gives a path of such edges from the
identity edge to its own edge. Conversely the inclusions of the two
factors make this the usual tree of $P*Q$ embedded in the ambient
tree. Denote the translate associated to the coset $gH$ by $T(gH)$.
Its edge labels are exactly the elements of $gH$.

Take finitely many pairwise intersecting cosets $g_jH_j$. Their
subtrees meet pairwise in edges and hence have a common vertex $v$,
by the finite Helly property for subtrees of a tree. Suppose
$v=xA$; the other case is identical. Each set
$g_jH_j\cap xA$ is nonempty and, after multiplication by $x^{-1}$,
is a left coset of $P_j$ in $A$. For two of our subtrees, choose an
edge whose label belongs to both original cosets. The path from that
edge to $v$ lies in both subtrees. Its edge incident with $v$ has a
single label belonging to both cosets and to $xA$. Thus these local
$P_j$ cosets meet pairwise. The assumed Helly property in $A$ gives a
common element, which belongs to every original coset. The tree
Helly fact follows, for example, by rooting the tree and choosing the
deepest of the nearest-to-root points of the finitely many subtrees.
\end{proof}

\begin{theorem}\label{thm:free}
Let $\alpha=(A,(A_i)_{i\in I_A})$ and
$\beta=(B,(B_i)_{i\in I_B})$ be coset incidence systems with disjoint
type sets. In $G=A*B$ put $G_i=A_i*B$ for $i\in I_A$ and
$G_i=A*B_i$ for $i\in I_B$. Then
\[
 G_J=A_{J\cap I_A}*B_{J\cap I_B}.
\]
The resulting system is a flag-transitive geometry if and only if
both input systems are flag-transitive. Under these conditions it is
residually connected if and only if both inputs are.
\end{theorem}
\begin{proof}
The displayed formula and forward flag assertion are the two preceding
lemmas, applied to the families of all input parabolics. For necessity,
lift a flag of $\alpha$ to the corresponding $G_i$ cosets using the
same representatives in $A$. Pairwise intersections remain nonempty.
A common element in $G$, followed by the retraction $G\to A$ that
kills $B$, gives a common element of the original cosets, since
$G_i$ maps onto $A_i$ for $i\in I_A$. This proves flag transitivity
of $\alpha$ without any Borel assumption. The retraction to $B$
gives the other input.

For a flag-transitive geometry, connectedness of a residue of rank at
least two is equivalent to generation of its group by the maximal
parabolics in that residue. To see this criterion, chambers incident
along an element connect their representatives through the appropriate
parabolic coset; the component containing the standard chamber consists
exactly of the chambers represented by the generated subgroup. Flag
transitivity identifies every residue with its parabolic coset system.

Apply this criterion to a face of type $J$. If both
$I_A\setminus J$ and $I_B\setminus J$ are nonempty, a remaining
$A$ type has residue parabolic containing the whole factor
$B_{J\cap I_B}$, and a remaining $B$ type has one containing
$A_{J\cap I_A}$. These two factors generate $G_J$. If only $A$
types remain, there are at least two of them. The residue parabolics
generate $G_J=A_{J\cap I_A}*B_{I_B}$ exactly when their projections
generate $A_{J\cap I_A}$, which is the required residue condition
in $\alpha$. The case with only $B$ types is symmetric. Taking
$J=J_A\cup I_B$, or $J=I_A\cup J_B$, proves necessity for every
input residue as well.
\end{proof}

\begin{proposition}\label{prop:panel}
Put $H_A=A_{I_A}$ and $H_B=B_{I_B}$. For a type $i\in I_A$,
write $U=A_{I_A\setminus\{i\}}$ and $U_0=H_A$.
The panel size in the free-product geometry is
\[
 [U*H_B:U_0*H_B]=
 \begin{cases}
 1,&U=U_0,\\
 [U:U_0],&H_B=1,\\
 \infty,&U_0<U\text{ and }H_B\ne1.
 \end{cases}
\]
There is the symmetric formula for a type in $I_B$. Consequently the
free-product geometry is firm if and only if both inputs are firm.
If both type sets are nonempty, it is thin if and only if both inputs
are thin and $H_A=H_B=1$.
\end{proposition}
\begin{proof}
The Borel is $H_A*H_B$ by Theorem~\ref{thm:free}, giving the index.
The first two cases follow directly. For the last, choose
$u\in U\setminus U_0$ and $1\ne v\in H_B$. The cosets
$(uv)^n(U_0*H_B)$, $n\geq0$, are distinct: equality would put a
nonzero power of $uv$ in $U_0*H_B$, contrary to its reduced word.
This proves infinitude, including when $U$ or $H_B$ is infinite.
An input panel has size one precisely when $U=U_0$, which proves
firmness. Under thinness all input panels must be nontrivial. If
both type sets occur, the last case then forces both Borels to be
trivial; the second case recovers exactly the input panel sizes.
The converse is immediate from the same formulas.
\end{proof}

The trivial-Borel condition cannot be omitted. Take both inputs to
be the rank-one geometry $(C_2^2,(C_2))$. Each has exactly two
vertices and is thin and flag-transitive. Their free-product geometry
has two types, and every panel is infinite by
Proposition~\ref{prop:panel}. This is a counterexample to the
unrestricted thinness assertion in Piedade--Tranchida,
arXiv:2505.24662v1, Corollary 3.9. Their definition permits
nonfaithful actions; imposing a regular action with trivial Borel
recovers the asserted thin case. Theorem~\ref{thm:free} also removes
the extra Borel assumption from their converse flag assertion.

\section*{Converse for a compatible amalgam}
We now give the converse missing from the separate tree gluing note.
Let $I_A,I_B$ be arbitrary finite type sets, put
$L=I_A\cap I_B$, $P=I_A\setminus L$, $Q=I_B\setminus L$,
and suppose a fixed isomorphism identifies $C=A_P$ with $B_Q$
and identifies $A_{P\cup J}$ with $B_{Q\cup J}$ for every
$J\subseteq L$. No flag assumption is imposed initially.
Set $G=A*_C B$ and
$G_i=\langle A_{\{i\}\cap I_A},B_{\{i\}\cap I_B}\rangle$.

\begin{proposition}\label{prop:converse}
If $(G,(G_i))$ is flag-transitive, both input incidence systems are
flag-transitive. Together with the tree gluing theorem, this is an
if and only if statement.
\end{proposition}
\begin{proof}
For $q\in Q$, we have $G_q=\langle A,B_q\rangle$ and
$C\leq B_q$. Fix left transversals for $C$ in $A$ and $B$.
A reduced amalgam word belongs to $\langle A,B_q\rangle$
precisely when all its $B$ syllables belong to $B_q$. This criterion
is independent of the choices of representatives, since multiplying
a syllable by an element of $C$ preserves membership in $B_q$.
It follows from uniqueness of the amalgam normal form, or by choosing
the $B$ transversal to extend one for $C$ in $B_q$.
Intersecting over $q\in Q$, every such syllable must belong to
$\bigcap_{q\in Q}B_q=C$, so there are no nontrivial $B$
syllables. Hence $G_Q=A$. If $Q$ is empty then $C=B$ and
$G=A$, so the same equality holds. Symmetrically $G_P=B$.

For $i\in I_A$, the two generating subgroups of $G_i$ have the
same intersection with $C$, by the fixed compatibility assumption.
The reduced word theorem embeds their amalgam in $G$ and gives
$G_i\cap A=A_i$. This is the single-subgroup embedding statement;
it does not use the false unrestricted intersection formula for two
special amalgam subgroups. In the flag-transitive geometry of $G$,
the residue of its standard $Q$ face is therefore the coset geometry
with group $G_Q=A$ and parabolics $G_Q\cap G_i=A_i$,
$i\in I_A$. Every residue of a flag-transitive geometry is
flag-transitive. This proves the assertion for $\alpha$, and the
standard $P$ residue gives $\beta$. Empty private type sets cause
no exception: the corresponding residue is the whole geometry.
\end{proof}

The reduced word theorems and trees are classical inputs. The explicit
Borel-dependent panel formula, the thinness counterexample and the
normal-form proof of the converse are the additions in this note.
They do not assert historical priority. The source's false general
special-subgroup intersection lemma remains corrected separately.
No full source recovery, numerical factor or formalization is claimed
by this private draft alone.
\end{document}
